% sample_chapter_answers.tex -- answers for the design-system sample chapter.
% In a real book this file is chN_answers.tex; main.tex puts it in the
\begin{chapteranswers}{2}{Uniform circular motion}
  \item One turn is the circumference, so
        $v = \dfrac{2\pi r}{T} = \dfrac{2\pi \times 50}{40} = \qty{7.85}{m.s^{-1}}$.
        Answer: \final{$v = \qty{7.9}{m.s^{-1}}$}.
  \item The car's speed stays the same, but its direction keeps changing. So its
        velocity changes, and a change in velocity is an acceleration. The
        acceleration points to the centre of the roundabout. Friction between
        the tyres and the road gives the force towards the centre.
  \item $F_\mathrm{c} = \dfrac{mv^2}{r} = \dfrac{0.50 \times 4.0^2}{1.2}
        = \qty{6.67}{N}$. Answer: \final{$F_\mathrm{c} = \qty{6.7}{N}$} towards
        the centre.
  \item $F_\mathrm{c}$ depends on $v^2$. Doubling $v$ multiplies $v^2$ by
        $2^2 = 4$. The force becomes 4 times bigger:
        $4 \times \qty{6.67}{N} = \qty{27}{N}$.
  \item $\omega = \dfrac{\Delta\theta}{t} = \dfrac{3\pi}{1.5} = 2\pi
        = \qty{6.28}{rad.s^{-1}}$. Then
        $v = \omega r = 6.28 \times 0.30 = \qty{1.88}{m.s^{-1}}$.
        Answers: \final{$\omega = \qty{6.3}{rad.s^{-1}}$} and
        \final{$v = \qty{1.9}{m.s^{-1}}$}.
\end{chapteranswers}

\begin{practiceanswers}[Practice question answers (sample)]
\begin{practiceans}
  \correct{C}
  The direction changes, so the velocity is not constant (A is wrong). A
  changing velocity needs a net force (B is wrong). There is no outward force
  (D is wrong). The acceleration points to the centre.
\end{practiceans}

\begin{practiceans}
  $v = 54 \div 3.6 = \qty{15}{m.s^{-1}}$. Friction gives the centripetal force:
  \[ F_\mathrm{c} = \frac{mv^2}{r} = \frac{1200 \times 15^2}{60} = \qty{4500}{N} \]
  Answer: \final{$F = \qty{4.5e3}{N}$} towards the centre of the curve.
  \begin{markingcriteria}
    \criterion{Changes \qty{54}{km.h^{-1}} to \qty{15}{m.s^{-1}}}{1}
    \criterion{Uses $F_\mathrm{c} = \dfrac{mv^2}{r}$ with correct substitution}{1}
    \criterion{Correct answer with unit (\qty{4.5e3}{N})}{1}
  \end{markingcriteria}
\end{practiceans}

\begin{practiceans}
  \begin{modelanswer}
  To go round a bend, a car needs a net force towards the centre,
  $F_\mathrm{c} = \dfrac{mv^2}{r}$. On a flat road, friction between the tyres
  and the road gives this force. Ice makes the largest possible friction force
  much smaller. At the same speed the car still needs the same $F_\mathrm{c}$,
  but friction can no longer supply it. Without enough inward force the car
  cannot follow the curve. It keeps going in a straighter line (Newton's first
  law), so it slides to the outside of the bend.
  \end{modelanswer}
  \begin{markingcriteria}
    \criterion{Explains that friction provides the centripetal force and that
      ice reduces the maximum friction, so the needed $F_\mathrm{c}$ cannot be
      supplied and the car moves to the outside of the bend}{3}
    \criterion{States that friction provides the centripetal force and that ice
      reduces friction}{2}
    \criterion{States one relevant point (for example, friction gives the force
      towards the centre)}{1}
  \end{markingcriteria}
\end{practiceans}
\end{practiceanswers}
